\begin{papercard}{nissenbaum2004}
\cf{Stages} \stagetag{S} (definition of a privacy violation as breach of context norms), \stagetag{H} (values and harms that decide whether a breach matters), \stagetag{C} (dimensions for coding an information flow). Nissenbaum, ``Privacy as Contextual Integrity'' (\emph{Washington Law Review} 79(1):119, 2004; PDF with footnotes). The canonical article-length statement of contextual integrity (CI), added to the corpus because the 2009 book is only available as contents and introduction. It is a normative argument written about \emph{public surveillance} (online public records, consumer profiling, RFID tags), not a method and not about health apps; the CI formulation here (norms of appropriateness and of flow) is the earlier one, before the book's explicit actor--attribute--transmission-principle scheme.

\cf{Problem} Public surveillance is among the least understood privacy challenges; US privacy policy is fragmentary and rests on ambivalent intuitions about shopper cards, CCTV and biometrics. The article says that prominent theories fail on these cases and proposes CI as ``an alternative benchmark for privacy'': it ``ties adequate protection for privacy to norms of specific contexts, demanding that information gathering and dissemination be appropriate to that context and obey the governing norms of distribution within it'' \pc{nissenbaum2004}{2}{Contextual integrity ties adequate protection for privacy to norms of specific contexts, demanding that information gathering and dissemination be appropriate to that context and obey the governing norms of distribution within it}. The claim is that public surveillance ``violates a right to privacy because it violates contextual integrity'' \pc{nissenbaum2004}{2}{public surveillance violates a right to privacy because it violates contextual integrity}.

\cf{Three motivating cases} (1) \emph{Public records online} \pc{nissenbaum2004}{3}{Case 1: Public Records Online}: records already public become a concern when put on the web, which looks paradoxical. (2) \emph{Consumer profiling and data mining}: people react with indignation although ``often the information in question is not confidential or sensitive in nature'' \pc{nissenbaum2004}{4}{Why? Often the information in question is not confidential or sensitive in nature}. (3) \emph{RFID tags} that could disseminate information about consumers without permission or awareness \pc{nissenbaum2004}{4}{Case 3: Radio Frequency Identification (RFID) Tags}. The cases show why a sensitive/non-sensitive, private/public split is inadequate.

\cf{The core idea} ``A central tenet of contextual integrity is that there are no arenas of life not governed by norms of information flow'' \pc{nissenbaum2004}{20}{A central tenet of contextual integrity is that there are no arenas of life not governed by norms of information flow}. Contexts range from broad spheres (education, politics, marketplace) to fine routines (visiting the dentist) and are partly constituted by norms from history, culture, law and convention. Two types of informational norms: \emph{appropriateness} and \emph{flow or distribution} \pc{nissenbaum2004}{21}{I posit two types of informational norms: norms of appropriateness, and norms of flow or distribution}. ``Contextual integrity is maintained when both types of norms are upheld, and it is violated when either of the norms is violated'' \pc{nissenbaum2004}{21}{Contextual integrity is maintained when both types of norms are upheld, and it is violated when either of the norms is violated}; a privacy complaint is sound when one of the norms has been transgressed \pc{nissenbaum2004}{21}{a complaint that privacy has been violated is sound in the event that one or the other types of the informational norms has been transgressed}, though a violation may be justified by a more urgent value \pc{nissenbaum2004}{21}{It still holds that a violation can be justified in the event that another, more serious or urgent value is at stake}.

\cf{Norms of appropriateness} They ``dictate what information about persons is appropriate, or fitting, to reveal in a particular context'' \pc{nissenbaum2004}{21}{norms of appropriateness dictate what information about persons is appropriate, or fitting, to reveal in a particular context}. Medical example: the patient shares physical-condition details with the physician but not vice versa \pc{nissenbaum2004}{21}{In medical contexts, it is appropriate to share details of our physical condition}. Norms vary in how explicit and complete they are (open in friendship, strict in a courtroom), and ``there is no place not governed by at least some informational norms'' \pc{nissenbaum2004}{22}{there is no place not governed by at least some informational norms}; the idea that nothing applies in public is ``pure fiction'' \pc{nissenbaum2004}{22}{no norms are in operation, that anything goes, is pure fiction}. Borrowed from Rachels and Schoeman: moving information from one situation into another can itself be a violation.

\cf{Norms of flow} Drawing on Walzer's theory of spheres of justice, ``another set of norms govern what I will call flow or distribution of information'' \pc{nissenbaum2004}{23}{another set of norms govern what I will call flow or distribution of information}. What matters is not only whether information is appropriate but ``whether its distribution, or flow, respects contextual norms of information flow'' \pc{nissenbaum2004}{24}{What matters is not only whether information is appropriate or inappropriate for a given context, but whether its distribution, or flow, respects contextual norms of information flow}. Flow principles named: free choice, discretion, confidentiality, and also need, entitlement and obligation, ``a list that is probably open-ended'' \pc{nissenbaum2004}{25}{Others include need, entitlement, and obligation}. In health care the flow is not governed by the patient's discretion but is close to mandated by the physician who may condition treatment on sharing, and it is not normally bidirectional; confidentiality of patient information is the subject of complex norms, with HIPAA cited \pc{nissenbaum2004}{25}{Another difference from friendship is that in the healthcare context, the flow is not normally bidirectional}.

\cf{Violation, conservatism and the presumption} A privacy violation occurs ``when either contextual norms of appropriateness or norms of flow have been breached'' \pc{nissenbaum2004}{26}{a normative account of privacy in terms of contextual integrity asserts that a privacy violation has occurred when either contextual norms of appropriateness or norms of flow have been breached}; information is always tagged with its context and never ``up for grabs'' because it is public \pc{nissenbaum2004}{26}{is always tagged with that context and never up for grabs as other accounts would have us believe}. Two worries: CI is ``conservative in possibly detrimental ways'' (e.g., looking up a surgeon's malpractice record online departs from past practice) \pc{nissenbaum2004}{26}{contextual integrity is conservative in possibly detrimental ways}, and, being tied to convention, it may lose moral authority \pc{nissenbaum2004}{27}{contextual integrity, being so tied to practice and convention, loses prescriptive value or moral authority}. The answer: ``a presumption in favor of the status quo'' \pc{nissenbaum2004}{28}{I propose that the requirement of contextual integrity sets up a presumption in favor of the status quo}, rebuttable by comparing entrenched norms with the novel practice in terms of context-internal goods \emph{and} fundamental social, political and moral values \pc{nissenbaum2004}{29}{entrenched norms be compared with novel practices that breach or threaten them, and judged worth preserving, or not, in terms of how well they promote not only values and goods internal to a given context, but also fundamental social, political, and moral values}.

\cf{Values used to evaluate breaches} (1) prevention of information-based harm, (2) informational inequality, (3) autonomy, (4) freedom, (5) preservation of important human relationships, (6) democracy and other social values \pc{nissenbaum2004}{29}{(2) informational inequality, (3) autonomy, (4) freedom, (5) preservation of important human relationships, and (6) democracy and other social values}; competing values supporting free flow are free speech, wealth, efficiency and security. \emph{Information-based harm}: unrestricted access ``can be harmful'' \pc{nissenbaum2004}{30}{Information in the wrong hands or generally unrestricted access to information can be harmful}, including identity theft and risks to employment, life and medical insurance \pc{nissenbaum2004}{30}{various goods such as employment, life, and medical insurance, could be placed at risk if the flow of medical information were not restricted}. \emph{Informational inequality}: in routine transactions people ``do not know what will be done with the information; do not grasp the full implications of consenting'' and cannot retract \pc{nissenbaum2004}{31}{individuals have little knowledge and understanding of the potential value of this economic exchange; do not know what will be done with the information; do not grasp the full implications of consenting to release of information}. A public-health argument adapted from Goldman: stronger protection of medical information makes people more likely to seek care and join research \pc{nissenbaum2004}{33}{individuals would then be more likely both to seek medical care and agree to participate in medical research}.

\cf{Applying CI to the cases} Public records online: the move from local to global accessibility ``constitutes a breach of entrenched norms of flow'' that requires a value assessment. Consumer profiling: ``the crucial issue is not whether the information is private or public \ldots\ but whether the action breaches contextual integrity'' \pc{nissenbaum2004}{35}{the crucial issue is not whether the information is private or public, gathered from private or public settings, but whether the action breaches contextual integrity}. A merchant using purchase records to stock goods and market to the same customer is within norms; a grocer who sells purchase data to magazine vendors and brokers breaches norms of appropriateness \emph{and} flow \pc{nissenbaum2004}{36}{The grocer who provides information about grocery purchases to vendors of magazine subscriptions or information brokers like Seisint and Axciom is responsible not only for breaches of norms of appropriateness but also norms of flow}. RFID: discretion is moved from the customer to information gatherers \pc{nissenbaum2004}{36}{Unless RFID tags are designed specifically to allow for easy detection and disabling, discretion is removed from the customer}.

\cf{Parameters to examine} The norms demand attention to ``the nature of the information in question and its relationship to the context, the roles involved in the context, the relationships among the roles, the rules of flow'' and effects on underlying values \pc{nissenbaum2004}{36}{the nature of the information in question and its relationship to the context, the roles involved in the context, the relationships among the roles, the rules of flow}. For third parties one should know ``their social roles, their capacity to affect the lives of data subjects, and their intentions'' and ask ``whether the information practice under consideration harms subjects; interferes with their self-determination; or amplifies undesirable inequalities'' \pc{nissenbaum2004}{38}{It is important to ask whether the information practice under consideration harms subjects; interferes with their self-determination; or amplifies undesirable inequalities in status, power, and wealth}. The conclusion lists the variables: context, nature of the information in relation to that context, roles of agents receiving the information, their relationships to subjects, on what terms the subject shared it, and the terms of further dissemination \pc{nissenbaum2004}{38}{several variables, including the nature of the situation, or context; the nature of the information in relation to that context; the roles of agents receiving information; their relationships to information subjects; on what terms the information is shared by the subject; and the terms of further dissemination}. Building a conclusive argument takes ``painstaking analysis of details'' \pc{nissenbaum2004}{37}{terms of contextual integrity involves painstaking analysis of details}, so protecting privacy ``will be a messy task'' \pc{nissenbaum2004}{39}{protecting privacy will be a messy task, requiring a grasp of concepts and social institutions as well as knowledge of facts of the matter}. Law is not the only means of preserving integrity; decency, etiquette and convention also regulate flows \pc{nissenbaum2004}{40}{Policy and law are not the only means of preserving contextual integrity}.

\cf{Limits (card)} A normative argument developed on three US public-surveillance cases; no empirical test and no coding procedure; the analysis of each case is a sketch (the author says only sketches were possible); the parameters are listed but not formalised as a fixed tuple; baseline norms are assumed to be knowable from practice, while the article admits practice and prescription can diverge; the value list is a synthesis of other scholars, not a harm taxonomy; US law and examples dominate; health-care is an illustrative context, not analysed.

\cf{Use for this thesis} (1) Cite this article for the definition of a violation (breach of appropriateness or flow norms of a context) and for the variable list that a coder can record for each flow in a policy: context, information type, sender/recipient roles and their relationship, terms under which the subject shared it, terms of onward flow. (2) The six values give a bridge from norm breach to harm: information-based harm, informational inequality (not knowing what happens to data or consequences of consent), autonomy, relationships; the harm tree can use them as intermediate nodes between a flow and Citron and Solove's harm types. (3) The health-care passage supplies the baseline for a health app: flows to the clinician are expected and shaped by treatment need, while flows to advertisers or brokers are not; an app may sit between care and consumer contexts, and which baseline applies must be stated per app. (4) State in the Decision Log which CI formulation is used (2004: appropriateness and flow; 2009: parameters of subject, sender, recipient, information type, transmission principle) and keep terms consistent. (5) The presumption in favour of the status quo means the thesis should describe the established norm it measures against, not assume it.
\end{papercard}
